\documentclass[aspectratio=169,10pt,spanish]{beamer}

\input{../../../_preambulo_beamer.tex}
\input{../../../_comandos_beamer.tex}

\selectlanguage{spanish}

\title{MA1001B - Modelación Estadística para la Toma de Decisiones}
\subtitle{Lección 38: Muestras pareadas (t antes-después)}
\author{}
\institute{Tecnol\'ogico de Monterrey}
\date{}

\begin{document}

\begin{frame}[plain]
  \vspace{1.2cm}
  {\Large\bfseries MA1001B - Modelación Estadística para la Toma de Decisiones\par}
  \vspace{0.7cm}
  {\large Lección 38: Muestras pareadas (t antes-después)\par}
  \vspace{0.7cm}
  {\color{TecRojo}\rule{\textwidth}{1.2pt}}
  \vspace{0.8cm}

  Facultad MA1001B\\[0.3cm]
  Periodo\\[0.3cm]
  Tecnol\'ogico de Monterrey
\end{frame}

\begin{frame}{La pregunta antes-después}
  \begin{alertblock}{Pregunta guía}
    ¿Cómo se prueba un cambio de layout cuando las mismas 25 estaciones se
    miden antes y después?
  \end{alertblock}

  \vspace{0.6em}
  Al final de esta lección, usted puede:
  \begin{itemize}
    \item formar diferencias pareadas $d=\text{antes}-\text{después}$;
    \item ejecutar una prueba $t$ de una muestra de $H_0:\mu_d=0$;
    \item contrastar el EE pareado con un EE de dos muestras no pareado;
    \item explicar por qué el análisis no pareado ignora el emparejamiento.
  \end{itemize}

  \vfill
  \scriptsize Todos los tiempos de manejo usados hoy son sintéticos. No se usan datos reales de empresa.
\end{frame}

\begin{frame}{La unidad es la estación, no la columna}
  \small
  Una celda de empaque sintética registra el tiempo de manejo en las mismas
  25 estaciones antes y después de un cambio de layout. La unidad
  observacional es la estación.

  \[
    d_i = \text{antes}_i - \text{después}_i,\qquad i=1,\ldots,25.
  \]

  \begin{exampleblock}{Hipótesis pareadas}
    $H_0:\mu_d=0$ frente a $H_1:\mu_d\neq 0$, con $\alpha=0.05$. Esta es una
    $t$ de una muestra sobre las 25 diferencias, no una $t$ de dos muestras
    sobre 50 tiempos.
  \end{exampleblock}
\end{frame}

\begin{frame}[fragile]{Paso 1: Diferencias emparejadas}
  \begin{columns}[T]
    \begin{column}{0.42\textwidth}
      \small
      \begin{lstlisting}
d = before - after
dbar = d.mean()
s_d = d.std(ddof=1)
      \end{lstlisting}

      \begin{block}{Salida verificada}
        $\bar{d}=2.720813$\\
        $s_d=1.525772$\\
        $r=0.973356$
      \end{block}
    \end{column}
    \begin{column}{0.58\textwidth}
      \centering
      \includegraphics[width=\textwidth]{../figuras/t_pareada_01_diferencias.png}
      \vspace{0.2em}
      \scriptsize Líneas emparejadas antes-después del Paso 1
    \end{column}
  \end{columns}
\end{frame}

\begin{frame}{La $t$ pareada es $t$ de una muestra sobre $d$}
  \small
  \[
    t=\frac{\bar{d}}{s_d/\sqrt{n}}=\frac{2.720813}{0.305154}=8.916186,
    \qquad gl=24.
  \]
  \[
    p=2P(T_{24}\ge 8.916186)=4.395592\times 10^{-9}.
  \]
  Como $|t|>t^*=2.063899$ y $p<\alpha$, la prueba pareada rechaza
  $H_0:\mu_d=0$.
\end{frame}

\begin{frame}[fragile]{Paso 2: Decisión pareada}
  \begin{columns}[T]
    \begin{column}{0.40\textwidth}
      \small
      \begin{lstlisting}
se = s_d / sqrt(n)
t = dbar / se
p = 2 * t_dist.sf(abs(t), 24)
      \end{lstlisting}

      \begin{block}{Salida verificada}
        $EE=0.305154$, $t=8.916186$\\
        $p=4.395592\times 10^{-9}$\\
        Decisión: rechazar $H_0$
      \end{block}
    \end{column}
    \begin{column}{0.60\textwidth}
      \centering
      \includegraphics[width=\textwidth]{../figuras/t_pareada_02_prueba_pareada.png}
      \vspace{0.2em}
      \scriptsize Histograma de diferencias del Paso 2
    \end{column}
  \end{columns}
\end{frame}

\begin{frame}{Paso 3: El análisis no pareado ignora el emparejamiento}
  \begin{columns}[T]
    \begin{column}{0.42\textwidth}
      \small
      EE pareado $=0.305154$, $p=4.395592\times 10^{-9}$, rechazar.

      \vspace{0.4em}
      EE de Welch no pareado $=1.847143$, $p=0.147291$, no rechazar.

      \vspace{0.4em}
      Correlación antes-después: $0.973356$.

      \vspace{0.5em}
      \textbf{Límite:} tratar las dos columnas como muestras independientes
      descarta el pareo y puede invertir la decisión con $\alpha=0.05$.
    \end{column}
    \begin{column}{0.58\textwidth}
      \centering
      \includegraphics[width=\textwidth]{../figuras/t_pareada_03_limite_no_pareada.png}
      \vspace{0.2em}
      \scriptsize Valores p pareado frente a no pareado del Paso 3
    \end{column}
  \end{columns}
\end{frame}

\begin{frame}{Verificación de conceptos}
  \begin{enumerate}
    \item ¿Por qué $d_i=\text{antes}_i-\text{después}_i$ es la unidad
      observacional correcta aquí?
    \item Calcule $EE=s_d/\sqrt{n}$ a partir de $s_d=1.525772$ y $n=25$.
    \item Si $p=4.395592\times 10^{-9}$ y $\alpha=0.05$, ¿cuál es la decisión
      pareada?
    \item ¿Por qué el $p=0.147291$ no pareado no refuta el hallazgo pareado?
  \end{enumerate}

  \vspace{0.6em}
  \begin{block}{Conexión con la Lección 39}
    La sesión puede continuar directamente con dos proporciones independientes
    en una prueba A/B, donde las muestras no están pareadas.
  \end{block}

  \vfill
  \scriptsize Fuente: OpenStax, \textit{Introductory Business Statistics 2e},
  Capítulo 10, Sección 10.6. CC BY-NC-SA.
\end{frame}

\appendix



\end{document}
